\chapter{Toolchain and Versions}
\label{app:toolchain}

This is the reproducibility baseline for the printed requests and diagnostics,
not a recommendation to keep an old patch release in production. Use an
isolated installation for the exercises. Apply supported security updates to
deployed services and rerun their compatibility tests.

\section{The Version Table}

\begin{center}
\small
\begin{tabular}{p{0.49\linewidth}p{0.38\linewidth}}
\toprule
Component & Verified baseline \tabularnewline
\midrule
.NET SDK / .NET and ASP.NET Core runtime & 10.0.303 / 10.0.11 \tabularnewline
Main target framework & \code{net10.0} \tabularnewline
All Hot Chocolate packages, including federation and authorization & 16.6.1 \tabularnewline
EF Core SQLite / JWT bearer package & 10.0.11 / 10.0.11 \tabularnewline
Node.js / npm & 24.15.0 / 11.13.0 \tabularnewline
Cosmo CLI (\code{wgc}) & 0.129.9 \tabularnewline
Embedded composition package & 0.63.3 \tabularnewline
Cosmo Router / its Go compiler & 0.341.0 / go1.26.6 \tabularnewline
Federation vocabulary emitted by default & v2.6 \tabularnewline
Compatibility: Hot Chocolate / target & 14.3.1 / \code{net8.0} \tabularnewline
Compatibility: EF Core / runtime & 9.0.19 / 8.0.31 \tabularnewline
\bottomrule
\end{tabular}
\end{center}

The compatibility example uses the same SDK to build its older target.
Go and the composition library are bundled components, not separate
installations. The federation package can emit through v2.7, but this example
uses its v2.6 default. Newer documentation does not establish what this
combination accepts.

Hot Chocolate 14 no longer receives security fixes. The ChilliCream security
policy checked on September 14, 2026 lists 16.x as receiving fixes and versions
below 16 as not receiving them~\autocite{hcsecurityrepair}. Appendix~\ref{app:hc14}
helps read and migrate older code, not start a deployment on that release.
Current documentation has moved beyond the book's Hot Chocolate patch; use
the printed source when reproducing an API.

\section{Install and Check}

Install the SDK, not only the runtime, from Microsoft's official download
page. Select the archived SDK above and the binary matching the operating
system and architecture~\autocite{dotnetdownloadrepair}. A portable
installation keeps the exercise independent of a newer SDK on the machine.
Check from the shell that will run the example:

\begin{minted}{text}
dotnet --version
dotnet --list-runtimes
node --version
npm --version
npm install -g wgc@0.129.9
wgc --version
router -version
\end{minted}

The router executable comes from the matching Cosmo release asset described
in chapter~\ref{ch:enter-the-router}, not from npm. Extract the platform's
archive, check its published checksum and make the executable available as
\code{router} in the exercise shell. On Windows its filename is
\code{router.exe}. Installing the CLI does not install it.

Keep all Hot Chocolate references within a service on the same version.
EF Core and JWT bearer follow .NET's release family, not Hot Chocolate's.
Build before exporting: a stale executable can export a valid schema that
does not describe the source just edited.

\section{Directories and Ports}

Appendix~\ref{app:source} prints each application file under its destination
path from this example root:

\begin{minted}{text}
src/
  Sessions/    localhost:5001   conference.db
  Speakers/    localhost:5002   speakers.db
  Ratings/     localhost:5003   ratings.db
  Search/      localhost:5004   search.db
graph/         graph.yaml and generated router.json
router/        config.yaml; router listens on localhost:3002
\end{minted}

Do not run the single-service compatibility example beside the main Sessions
service: both use port 5001. From each service directory, run:

\begin{minted}{text}
dotnet restore
dotnet build
dotnet run --no-build -- schema export --output schema.graphql
dotnet run --no-build
\end{minted}

The last command stays running in its own terminal. Repeat for the other
services needed by the current chapter. The supplied launch profile selects
the port and Development environment, which enables the statement log.

From the example root, compose and launch in another terminal:

\begin{minted}{text}
cd graph
wgc router compose -i graph.yaml -o router.json
cd ../router
router -config config.yaml
\end{minted}

Schema paths are relative to the graph input. The router's execution-config
path is relative to the router's working directory. After a schema change,
export, compose, then restart this statically configured router.

\section{Reset Only the Exercise Data}

The example uses \code{EnsureCreated}, not migrations. After changing a seed
or model, stop the affected service and remove only its disposable database,
then restart. In PowerShell, starting from the example root:

\begin{minted}{text}
cd src/Sessions
Remove-Item -LiteralPath ./conference.db
dotnet run
\end{minted}

Do not reset a database containing data that must be kept. A successful
mutation changes a stored row until reset; rebuilding does not restore the
seed. Chapters~\ref{ch:schema-design} and~\ref{ch:representation-order}
depend on the original start-time order.

The signing key and token in chapter~\ref{ch:entities-authorization} are
development fixtures. Keep the example on a trusted local machine. They do
not replace an identity provider, production token validation, protected
subgraph ports or a secret store.
